\subsection{Problem setup and localization diagnostics}

We consider ordinal classification with ordered labels
\(\mathcal{Y}=\{0,\ldots,K-1\}\). The upper endpoint \(K-1\) is the
rare-end class in the settings studied here. A model maps an input \(x\) to a
penultimate representation \(h(x)\in\mathrm{R}^d\), logits
\(z_k(x)=w_k^\top h(x)+b_k\), and class probabilities
\(p_k(x)=\exp\{z_k(x)\}/\sum_{j=0}^{K-1}\exp\{z_j(x)\}\). Our aim is not to
propose a new classifier or to establish that class imbalance is an
independently isolated cause of error. Instead, we characterize a recurring
empirical pattern in the studied imbalanced settings: predictions for a rare
upper endpoint are displaced toward interior ordinal classes. We call this
pattern \emph{rare-end inward localization}.

This terminology deliberately separates a model's probability distribution,
its discrete decision, and the geometry of its learned representation. The
modal decision is \(\hat y_{\mathrm{mode}}(x)=\arg\max_k p_k(x)\), whereas the
exact discrete action for absolute ordinal loss is
\[
\hat y_{L1}(x)=\arg\min_{a\in\{0,\ldots,K-1\}}
    \sum_{k=0}^{K-1}p_k(x)|k-a|.
\]
Thus \(\hat y_{L1}\) is a posterior median (with the deterministic tie rule
used in the saved evaluation code), rather than a heuristic rounding of the
predictive mean. We report both decisions when relevant and use the exact L1
action for the primary ordinal-error summaries. An error is severe when
\( |Y-\hat Y|\geq 2 \). This threshold is descriptive of errors that skip at
least one adjacent ordinal category; it is not a new uncertainty metric.

For a predicted distribution, its \emph{predictive location} is
\[\mu_p(x)=\sum_{k=0}^{K-1} k p_k(x).\]
For a true upper-endpoint observation (\(Y=K-1\)), we define \emph{inward
shrinkage} as \[S(x)=(K-1)-\mu_p(x).\]
Larger positive \(S\) means that probability mass is located farther inward
from the upper endpoint; it is not the same quantity as discrete MAE. We
therefore inspect routing counts, endpoint and adjacent-class probabilities,
predictive location, shrinkage, L1 MAE, and severe-error prevalence together.
The corresponding lower-endpoint diagnostics are retained as controls where
reported, because a change that improves one endpoint can have asymmetric
consequences elsewhere in the ordinal label space.

All rare-end quantities condition on the true endpoint population rather than
on samples predicted as the endpoint. This distinction is essential: a model
that almost never routes a true endpoint to \(K-1\) can otherwise make the
remaining predicted-endpoint subset appear deceptively easy. We summarize both
the full rare-end population and its discrete routing distribution. In
particular, endpoint probability \(p_{K-1}\), adjacent probability
\(p_{K-2}\), and their joint allocation indicate whether an apparent movement
in \(\mu_p\) is accompanied by greater endpoint support or simply a different
spread among interior classes. These probability summaries complement, but do
not replace, realized ordinal decisions. Conversely, a change in a decision
rule cannot be interpreted as a change in representation or in the original
probability distribution. This conditioning and decomposition are used
throughout to avoid conflating localization, calibration-like probability
behavior, and decision-dependent loss.

The two backbone objectives are CE and the Ranked Probability Score (RPS). For
a one-hot label \(y\), the per-example RPS used here is
\[
\operatorname{RPS}(p,y)=\frac{1}{K-1}\sum_{r=0}^{K-2}
\left(\sum_{k=0}^{r}p_k-\mathbf{1}\{y\leq r\}\right)^2.
\]
CE is the nominal probabilistic baseline. RPS evaluates squared errors in
cumulative probability and is an ordinal probabilistic baseline; the paper
does not treat it as universally superior. Both objectives train backbones,
while the head adaptation below uses a common controlled objective so that its
interpretation does not rest on changing the representation simultaneously.

Our representation diagnostics are training-derived and descriptive. For a
given fitting set, the class centroid is
\[\bar h_k=\frac{1}{n_k}\sum_{i:y_i=k} h(x_i).\]
For a rare-end point, the endpoint-versus-adjacent margin is
\[
m_{\mathrm{adj}}^{\mathrm{rep}}(x)=
\lVert h(x)-\bar h_{K-2}\rVert_2-
\lVert h(x)-\bar h_{K-1}\rVert_2.
\]
A larger value means that the endpoint is closer than the immediately inward
class. We also use the order-agnostic generic centroid gap
\(g^{\mathrm{centroid}}(x)=d_{(2)}(x)-d_{(1)}(x)\), where \(d_{(1)}\) and
\(d_{(2)}\) are the two smallest distances to any class centroids. These
quantities help distinguish endpoint-specific alignment from generic feature
separation. Nearest-centroid routing is a diagnostic summary only: it neither
defines the classifier nor establishes a causal representation mechanism.

The central intervention is designed around the decomposition
\[\text{representation}\ \longrightarrow\ \text{probabilistic head}
\ \longrightarrow\ \text{decision rule}\ \longrightarrow\ \text{ordinal risk}.\]
Write each class row as \(w_k=s_kv_k\), where \(s_k=\lVert w_k\rVert_2\) and
\(\lVert v_k\rVert_2=1\). Condition A is the original classifier head,
including its directions, scales, and biases, applied to a frozen learned
representation. Historical condition B is a balanced full-head refit, in which
a new linear head is trained with class-balanced sampling and all head
parameters may change. It is a comparative probe, not the primary
confirmatory contrast.

Condition C initializes from A but adapts only the unit directions:
\(w_k^C=s_k^A v_k^C\). Every original per-class scale \(s_k^A\) and bias
\(b_k^A\) is fixed. Consequently, an A--C difference is evidence about the
head-actionable portion of localization for that fixed representation, not
evidence that the representation itself was improved. C is a mechanistic probe,
not a proposed deployment method. The historical D controlled-scale probe
changes a prescribed scale component while retaining the direction-control
logic; it is not part of the primary confirmatory comparison and its detailed
mechanics are supplementary. Here, class-balanced sampling means
replacement-based resampling to form balanced batches; it is an established
control/intervention rather than a novel contribution.
